\documentclass[11pt,a4paper]{article}
\usepackage[margin=2.4cm]{geometry}
\usepackage{amsmath,amssymb,booktabs,graphicx,xcolor}
\usepackage{times}
\usepackage[colorlinks,linkcolor=blue!50!black,citecolor=blue!50!black]{hyperref}
\newcommand{\dd}{\,\mathrm{d}}
\newcommand{\sgn}{\operatorname{sgn}}
\newcommand{\TODO}[1]{\textcolor{red}{[#1]}}
\title{\bf Optimal-transport targets for macromolecular model fitting:\\
non-vanishing gradients from a structure-factor computation}
\author{P.\ H.\ Zwart\\
\small Molecular Biophysics and Integrated Bioimaging, Lawrence Berkeley National Laboratory}
\date{}
\begin{document}\maketitle

\begin{abstract}\noindent
Real-space refinement targets are local: their gradient decays with the overlap between
model and map, so a model placed even a few resolution widths away receives no usable
signal. Optimal-transport (OT) discrepancies do not share this behaviour --- the gradient
of a Wasserstein-1 objective is bounded away from zero wherever mass is misplaced --- but
the formulations used to date represent both densities as point clouds in real space, which
does not compose with refinement engines built on structure factors. We show that sliced
$W_1$ against a sampled map \emph{is} a structure-factor computation: the target's
projected profile along any direction is a central slice of its transform, the model's
projected profile has an exact one-dimensional transform which is a structure factor with
the usual atomic form factor, and convolution, differentiation and integration all become
reciprocal-space multipliers. Atomic gradients then follow by the route crystallography
already uses, with atoms gridded at a common width and the residual $B$ applied as a
per-type multiplier. Against a case with a closed-form answer the gradient agrees with
finite differences to $3\times10^{-9}$; the floor reaches $5\times10^{-7}$ on cubic,
non-cubic, monoclinic and triclinic grids alike; and the gridded route is $152\times$
faster than direct summation at $8\times10^{3}$ atoms. We argue that the resulting target
is a \emph{search} device to be followed by a conventional forward-model refinement, and
we are explicit that the protocols needed to make it robust in production are outside the
scope of this work.
\end{abstract}

\section{Introduction}

\subsection{Why refinement targets go blind}

A real-space refinement target scores a model by how its computed density agrees with an
observed map, and every such score is built from products of the two densities evaluated
where they overlap. When a model is misplaced by more than a resolution width that overlap
falls off like the tail of a Gaussian, and with it the gradient. The practical consequence
is familiar: a ligand modelled into the wrong lobe of density, or a side chain in the wrong
rotamer, is not repaired by refinement, because the target has nothing to say about where
the model should go. It has to be moved by something else --- a human, a search, or a
docking calculation --- and refinement then polishes whatever it is handed.

This is not a defect of any particular target so much as a property of local comparison.
A least-squares residual on the map, a correlation coefficient, and the sum of density at
atomic centres all share it, and they degrade in two distinguishable ways
(\S\ref{sec:reach}): the sum-of-density gradient keeps pointing roughly the right way but
shrinks by nine orders of magnitude over ten \AA ngstr\"om, while the least-squares
gradient keeps its magnitude and \emph{reverses direction} once the two densities cease to
overlap and the model's self-term dominates.

\subsection{Optimal transport, for crystallographers}

Optimal transport supplies a different way to compare two densities, and one that does not
require them to overlap at all. Given two non-negative densities of equal total mass, the
Kantorovich problem asks for the cheapest way to rearrange one into the other, where the
cost of moving a unit of mass from $x$ to $y$ is some ground cost $c(x,y)$. Taking
$c(x,y)=\lVert x-y\rVert$ gives the Wasserstein-1 distance $W_1$, sometimes called the
earth-mover's distance: the answer is literally how far mass has to travel, and it is
finite and informative however far apart the two densities are.

The reason this is not routine in crystallography is cost. Solving a transport problem
between two three-dimensional maps is a large linear program, and the entropically
regularised solvers that made OT tractable in machine learning introduce a regularisation
parameter whose setting is itself a problem. The device that makes OT affordable is
\emph{slicing}. In one dimension the transport problem has a closed-form solution obtained
by sorting: the optimal map is the monotone rearrangement, and
\begin{equation}
W_1(\mu,\nu)=\int \bigl|F_\mu(t)-F_\nu(t)\bigr|\dd t
\label{eq:w1d}
\end{equation}
where $F$ denotes a cumulative distribution. Sliced OT replaces the $d$-dimensional problem
by an average over one-dimensional problems obtained by projecting both measures onto a set
of directions,
\begin{equation}
\mathcal{SW}_1(\mu,\nu)=\frac{1}{L}\sum_{\ell=1}^{L}
 W_1\bigl(P_{u_\ell}\mu,\,P_{u_\ell}\nu\bigr),
\label{eq:sw1}
\end{equation}
each of which is a sort. The result is a genuine metric, it inherits the topology of $W_1$,
and it costs $O(N\log N)$ per direction rather than the cost of a linear program.

Sliced OT is by now well established in imaging and increasingly in cryo-EM. AlignOT
\TODO{cite Riahi et al., IEEE/ACM TCBB 20, 3842 (2023)} aligns two EM maps represented as
point clouds by minimising a Wasserstein objective over rotations, reporting convergence
from a wider range of starting angles than local alignment. EMPOT \TODO{cite Riahi et al.,
PRX Life 3, 023003 (2025)} extends this to partial overlap and rigid-body fitting of atomic
models using an unbalanced Gromov--Wasserstein divergence. Related work uses transport for
interpolation between maps \TODO{MorphOT}, for manifold learning over conformational
ensembles \TODO{Zelesko et al., ISBI 2020}, and for three-dimensional reconstruction
\TODO{CryoSWD, Zehni \& Zhao, ICASSP 2023}. Most directly, a recent study
\TODO{bioRxiv 2025.12.23.696001} reports that a sliced-Wasserstein image loss
``significantly broadens the basin of attraction'' relative to mean-squared error in
cryo-EM optimisation. \emph{That OT opens gradients is therefore established}, and it is
not what this paper claims. What is missing is a formulation that a crystallographic
refinement engine can actually execute.

\subsection{Non-vanishing gradients: the good and the bad}

The property that makes $W_1$ attractive follows from its dual. The optimal potential is
$1$-Lipschitz and, in one dimension, has the closed form
\begin{equation}
f'_\ell(t) = -\sgn\bigl(F_{P_{u_\ell}\mu}(t)-F_{P_{u_\ell}\nu}(t)\bigr),
\label{eq:pot}
\end{equation}
a \emph{sign}. Its magnitude is one wherever the two cumulative distributions differ at all,
and it is defined everywhere rather than only where the densities overlap. Smoothing by the
atomic form factor and backprojecting gives the atomic gradient
(\S\ref{sec:formulation}), whose magnitude is therefore bounded below as long as any mass
is misplaced, at any separation.

That is exactly the wanted behaviour for a search, and exactly the wrong behaviour for a
refinement. Because the gradient is a smoothed sign, it does not diminish as the model
approaches the answer, and it therefore carries no natural step length: a descent scheme
driven by it overshoots and limit-cycles rather than converging. Two consequences follow
and we adopt both. First, the step length must come from somewhere else; the barycentric
projection of the same transport plan --- the monotone rearrangement itself --- supplies a
displacement rather than a force, and it \emph{does} vanish at the solution
(\S\ref{sec:step}). Second, and more importantly, the transport term should not be the
final target at all. Left unrestrained it will happily distort a model to fit density: we
measure a configuration with $1.25$\,\AA{} distance errors scoring $8.5\times$ better on
the transport term than the chemically exact answer. The intended role is to deliver a
model into the correct basin, after which a conventional forward-model refinement --- with
its proper treatment of scaling, bulk solvent, and stereochemistry --- takes over.

\subsection{Why not just search harder}

The obvious alternative to a better target is a better search: exhaustive rotation--
translation scans, docking, or simulated annealing. These work, and for rigid-body placement
of a whole molecule they are standard. The difficulty is atomic-level fitting, where the
search space is not six-dimensional but $3N$-dimensional, and where the interesting failures
--- a wrong rotamer, a peptide flip, a ligand bound in a shifted register --- are not
reachable by rigid-body moves. Conformer-ensemble methods address this by enumerating
plausible geometries in advance, but the number of conformers grows quickly and each must
still be placed and scored. What is wanted is a target whose gradient reaches, so that a
small number of starting points suffices, and which is cheap enough to evaluate that an
ensemble of them can be scored at all. The formulations cited above are point-cloud based
and rigid-body; they do not provide per-atom analytic gradients, and they scale poorly in
the atom count where a protein fragment or a cofactor is concerned.

\subsection{What this paper does}

We show that sliced $W_1$ against a sampled map is a structure-factor computation, and that
its atomic gradients can therefore be evaluated by the FFT route crystallography has used
since Agarwal \TODO{cite Agarwal 1978} and Ten Eyck \TODO{cite Ten Eyck 1973/1977}: grid the
atoms at a common width, apply the residual $B$ as a per-type reciprocal-space multiplier,
transform, and gather at the atomic positions. The contribution is the formulation and its
cost, not an application. We demonstrate that naive, hand-set targets are already enough to
recover badly placed models in two synthetic cases, and we are explicit that the protocols
required to make this robust on experimental data --- masking and segmentation decisions,
per-atom $B$ treatment, handling of negative density, symmetry mates and the surrounding
structure --- are identified but not solved here.

\section{Formulation}\label{sec:formulation}

\TODO{Slice theorem; exact 1-D structure factor for the model; operators as multipliers;
the phase-origin and branch conventions. Text exists in the earlier draft --- port it.}

\section{Atomic gradients by the Agarwal/Ten Eyck route}\label{sec:fft}

\TODO{The $B$-factor split; forward scatter and backward gather; per-type grouping;
the observation that the gather kernel is the form factor itself rather than its
derivative, because the $2\pi i q$ from the phase cancels against the $1/2\pi i q$ of the
spectral antiderivative.}

\begin{table}[h]\centering\small
\caption{Cost of value and gradient, $L=24$, CPU, single precision-independent.
$K$ grows with the model because the projection window must contain it, so direct
summation scales worse than linearly in $N$.}
\begin{tabular}{@{}rrrrr@{}}\toprule
$N$ & $K$ & direct (ms) & gridded (ms) & speedup\\\midrule
23 & 95 & 2.6 & 3.9 & 0.7$\times$\\
200 & 119 & 22.0 & 5.6 & 4.0$\times$\\
2000 & 257 & 4594 & 64.7 & 71$\times$\\
8000 & 467 & 26079 & 171.4 & \textbf{152}$\times$\\\bottomrule
\end{tabular}
\end{table}

\section{Reach, and how conventional targets fail}\label{sec:reach}

Capped leucine fragment (13 heavy atoms, $N$-acetyl-leucine-$N'$-methylamide), $2.5$\,\AA,
$\sigma=1.062$\,\AA. ``cos'' is the cosine between $-\nabla E$ and the true displacement.

\begin{table}[h]\centering\small
\begin{tabular}{@{}rrrrrrrr@{}}\toprule
$t$ (\AA) & $t/\sigma$ & \multicolumn{2}{c}{OT} & \multicolumn{2}{c}{least squares}
 & \multicolumn{2}{c}{sum of density}\\
\cmidrule(lr){3-4}\cmidrule(lr){5-6}\cmidrule(lr){7-8}
 & & $\lVert g\rVert/w$ & cos & $\lVert g\rVert$ & cos & $\lVert g\rVert$ & cos\\\midrule
 2 & 1.9 & 0.5009 & 0.996 & 1.6e-5 & 0.988 & 1.6e-5 & 0.992\\
 4 & 3.8 & 0.5009 & 0.996 & 1.6e-5 & 0.976 & 2.8e-6 & 0.955\\
 6 & 5.7 & 0.5009 & 0.996 & 1.4e-5 & \textbf{0.431} & 5.6e-8 & 0.914\\
 8 & 7.5 & 0.5009 & 0.996 & 1.4e-5 & \textbf{$-$0.280} & 1.0e-10 & 0.895\\
12 & 11.3 & \textbf{0.5009} & \textbf{0.996} & 1.4e-5 & 0.229 & 1.6e-14 & $-$0.962\\
\bottomrule\end{tabular}
\end{table}

Radius of convergence, random rotation ($0$--$30^\circ$) plus translation of magnitude $d$,
10 trials, identical restraint projection and identical normalised step for every target,
no annealing, success at RMSD $<0.5$\,\AA:

\begin{table}[h]\centering\small
\begin{tabular}{@{}rrrrr@{}}\toprule
$d$ (\AA) & $d/\sigma$ & OT & least squares & sum of density\\\midrule
1--3 & $\le$2.8 & 100\% & 100\% & 100\%\\
4 & 3.8 & \textbf{100\%} & 100\% & 90\%\\
6 & 5.7 & \textbf{100\%} & 90\% & 80\%\\
8 & 7.5 & \textbf{100\%} & 10\% & 0\%\\
\bottomrule\end{tabular}
\end{table}

Normalising the step for every target deliberately favours the conventional ones, since it
neutralises their magnitude decay and tests only direction; the advantage reported here is
therefore a lower bound.

\section{Step length from the barycentric projection}\label{sec:step}

\TODO{$M^{-1}$ backprojection factor, one-step recovery of a rigid translation, and the
measured comparison: final RMSD $0.010$\,\AA{} with $\delta$-derived steps against
$0.133$\,\AA{} with a fixed cap, oscillation $0.002$ against $0.057$\,\AA.}

\section{Validation}\label{sec:valid}

\begin{table}[h]\centering\small
\caption{End-to-end validation across grid geometries, leucine fragment, $2.5$\,\AA,
$L=32$. The floor is a diagnostic of the \emph{map}: it detects inadequate padding, and
requires a clearance of $\gtrsim7\sigma$ between the model and the box or cell boundary.}
\begin{tabular}{@{}lrrrrr@{}}\toprule
grid & floor & anchor & FD median & reach @ 12\,\AA & 1-step\\\midrule
cubic $48^3$ & 4.79e-7 & $+0.000\%$ & 1.3e-9 & $+0.000\%$ & 5.2e-3\,\AA\\
non-cubic, anisotropic step & 4.76e-7 & $+0.000\%$ & 7.7e-10 & $+0.000\%$ & 5.2e-3\,\AA\\
monoclinic $60\times66\times54$ & 4.73e-7 & $+0.000\%$ & 1.7e-9 & $+0.000\%$ & 5.0e-3\,\AA\\
triclinic $96\times100\times104$ & 4.80e-7 & $+0.000\%$ & 2.7e-9 & $+0.000\%$ & 5.0e-3\,\AA\\
\bottomrule\end{tabular}
\end{table}

The formulation is exact on a general grid because the target slices are computed as a
structure factor over voxels rather than by interpolating a central slice of the map's
transform; the only place grid geometry enters is each voxel's Cartesian position.
Bandwidth is set by the shortest reciprocal vector of the sampling lattice, which on a
skewed cell can be twice what the step lengths suggest.

\section{Demonstrations}\label{sec:demo}

\TODO{NOT YET RUN. Both use synthetic data, equal $B$ factors, a naive hand-set stage
schedule, the alternating-projection protocol ($P_{\rm data}$ then $P_{\rm restr}$, over-
relaxed), and a final map-correlation refinement. In both cases the same protocol started
from a map-correlation target is the control.}

\subsection{Pentapeptide with wrong rotamers and backrub main-chain motion}
\TODO{Design: build a pentapeptide, perturb two side chains to incorrect rotamers and apply
a backrub motion to the intervening main chain; refine; report RMSD by residue and
separately for main chain and side chain.}

\subsection{NAD\textsuperscript{+}/NADH from a conformer ensemble}
\TODO{Design: generate a conformer ensemble, place each into the map of the correct
conformer, score by transport, refine the best; report rank of the correct conformer and
final RMSD.}

\section{Scope and what is deliberately left out}

\TODO{Search target not refinement target; balanced formulation and what breaks for a
fragment in a larger assembly; negative density and the segmentation decision (measured:
contouring truncates the second moment by up to 25\% and degrades the fit by 6$\times$);
per-atom $B$; symmetry mates and environment. Electron scattering factors are also Gaussian
sums, so the same route serves cryo-EM.}

\end{document}
